% Sign-in and institutional authentication: the two methods (password, proxy), their settings,
% proxy sign-in's behavior and security requirements, the setup procedure, example proxy
% configurations, accounts without a password, adding a method in code, the viewers' own sign-in.
% Sources: docs/sign-in.md; the code is the authority: medorama/src/signin.rs (Method, identify,
% login, logout, method), config.rs (sign_in_method), server.rs (session_user, my_password,
% user_add), users.rs (get_by_login, unusable_password), main.rs (user add --no-password),
% web/js/main.js (showLogin), web/js/views/admin.js; tests: medorama/tests/signin.rs.
\chapter{Sign-in and institutional authentication}
\label{chap:sign-in}

This chapter explains how users sign in to Medorama, and how an institution connects Medorama to
its own sign-in: single sign-on through OpenID Connect or SAML (for example Microsoft Entra ID,
Okta, Google or Shibboleth), Kerberos, multi-factor authentication or a VPN. It is written for
administrators who deploy Medorama and for integrators who adapt it. What users may see once
they are signed in is explained in \cref{chap:protecting}, and the security review is
\cref{chap:security}.

\section{Two ways to sign in}
\label{sec:sign-in-methods}

Medorama separates two questions:
\begin{itemize}
\item \textbf{Who is asking?} This is the job of the \emph{sign-in method}. All of it lives in
  one module, \file{medorama/src/signin.rs}.
\item \textbf{What may that user do and see?} Medorama decides this in the same way for every
  method, from the user's account: the role (viewer, analyst or administrator), the PHI
  permission and Privacy Mode (\cref{chap:users-privacy}).
\end{itemize}

\screenshot{sign-in.pdf}{Sign-in is one seam: either method identifies the user, and everything
after it (roles, the PHI permission, Privacy Mode, the privacy gate) is the same.}{fig:sign-in-seam}

The setting \code{sign\_in} in \file{medorama.toml} chooses the method
(\cref{tab:sign-in-methods}).

\begin{xltabular}{\linewidth}{@{}P{2.6cm}P{5.6cm}L@{}}
\caption{The two sign-in methods.}
\label{tab:sign-in-methods}\\
\toprule
\tabhead \code{sign\_in} & \tabhead Who the user is & \tabhead Typical use \\
\midrule
\endfirsthead
\toprule
\tabhead \code{sign\_in} & \tabhead Who the user is & \tabhead Typical use \\
\midrule
\endhead
\bottomrule
\endfoot
\code{password} (the default) & The account whose name and password were entered on Medorama's
  sign-in form & A single team, a test installation, a server reached through an SSH tunnel or a
  VPN \\
\code{proxy} & The account named in a request header by the institution's authenticating
  reverse proxy, on requests that come from that proxy & Single sign-on (OpenID Connect, SAML),
  Kerberos, multi-factor authentication, an access gateway \\
\end{xltabular}

With \code{proxy}, the proxy performs the institution's own sign-in, including any multi-factor
authentication, and Medorama never sees a password. Accounts, roles and the PHI permission still
live in Medorama, under the names the proxy gives the users. With either method, every request
handler receives the signed-in user in the same form (\code{server::Auth(User)}), so the rest of
Medorama does not depend on the method.

\section{Settings}
\label{sec:sign-in-settings}

The settings in \cref{tab:sign-in-settings} are read from \file{medorama.toml} in the catalog's
home folder (\file{config.rs}, \code{sign\_in\_method}). Medorama reads them when it starts, so
restart it after a change. Every setting is also listed in \cref{chap:settings-reference}.

\begin{xltabular}{\linewidth}{@{}P{3.4cm}P{3.0cm}L@{}}
\caption{The sign-in settings of \file{medorama.toml}.}
\label{tab:sign-in-settings}\\
\toprule
\tabhead Setting & \tabhead Default & \tabhead Meaning \\
\midrule
\endfirsthead
\toprule
\tabhead Setting & \tabhead Default & \tabhead Meaning \\
\midrule
\endhead
\bottomrule
\endfoot
\code{sign\_in} & \code{"password"} & \code{"password"} or \code{"proxy"}. \\
\code{proxy\_user\_header} & \code{"X-Remote-User"} & The request header in which the proxy
  sends the user's account name. Proxy mode only. \\
\code{proxy\_addresses} & \code{["127.0.0.1", "::1"]} & The IP addresses the proxy connects
  from. The header is believed only on requests from these addresses. Each entry is one
  address; ranges are not accepted. Proxy mode only. \\
\code{proxy\_logout\_url} & \code{""} & Where the page's \ui{Sign out} leads: the proxy's own
  sign-out address. When it is empty, the page shows its signed-out screen. Proxy mode only. \\
\code{secure\_cookies} & \code{false} & Marks Medorama's cookies \code{Secure}. Set it to
  \code{true} whenever users reach Medorama over HTTPS, with either method. Medorama warns at
  start when it listens on another address than this computer's, or signs users in through a
  proxy, with \code{secure\_cookies = false}. \\
\code{session\_hours} & \code{12} & In password mode, how long a sign-in lasts. In proxy mode,
  how long a visit's Privacy Mode choice is kept at most. \\
\code{session\_idle\_minutes} & \code{60} & In password mode, a session unused this long ends (0:
  only \code{session\_hours} ends it). In proxy mode, the gap after which a user's next request
  counts as a new visit in the audit log. \\
\end{xltabular}

Medorama refuses to start when these settings are inconsistent, so that a half-made change is
not silently ignored. \Cref{tab:sign-in-errors} lists the messages, which follow the path of the
settings file.

\begin{xltabular}{\linewidth}{@{}P{5.0cm}L@{}}
\caption{Start-up errors of the sign-in settings.}
\label{tab:sign-in-errors}\\
\toprule
\tabhead Mistake & \tabhead Message \\
\midrule
\endfirsthead
\toprule
\tabhead Mistake & \tabhead Message \\
\midrule
\endhead
\bottomrule
\endfoot
A \code{proxy\_\dots} setting without \code{sign\_in = "proxy"} & \code{proxy\_user\_header,
  proxy\_addresses and proxy\_logout\_url need sign\_in = "proxy"} \\
An unknown method & \code{sign\_in: "saml" is not a sign-in method (password, proxy)} \\
A header name with spaces or other invalid characters & \code{proxy\_user\_header: "bad header"
  is not a header name} \\
A host name, or a range, instead of an IP address & \code{proxy\_addresses: "proxy.example.org"
  is not an IP address} \\
An empty list of addresses & \code{proxy\_addresses: name the address the proxy connects
  from} \\
\end{xltabular}

The page asks the server which method is in use through the public request \code{GET
api/v1/sign-in}, which answers, for example, \code{\{"method": "proxy", "logout\_url":
"/oauth2/sign\_out"\}} (\cref{chap:http-api}).

\section{Sign-in with Medorama's passwords}
\label{sec:sign-in-password}

With \code{sign\_in = "password"}, users sign in on Medorama's own form.
\begin{itemize}
\item \textbf{Passwords} have at least 10 characters and are stored as Argon2id hashes (the
  \code{argon2} crate's defaults: 19\,MiB of memory, 2 iterations, 1 lane, and a random salt
  per password).
\item \textbf{Throttling.} Each attempt is counted before its password is checked, so attempts
  sent at once count too. After 5 failed attempts within 60 seconds, an account name is locked
  for 60 seconds from the address they came from; others, elsewhere, can still sign in under
  that name, so nobody can lock a known account out. After 30 failed attempts within 60 seconds
  from one address, whatever the names, that address is refused for 60 seconds. (Behind a
  reverse proxy, every request comes from the proxy's address: the limits are then the whole
  site's.) At most 4 passwords are checked at once. Every failure waits one second before the
  answer. A sign-in with an unknown name takes as long as one with a wrong password, so the
  timing does not reveal which names exist; a name no account can have (over 64 characters, or
  with other characters than letters, digits and \code{. \_ - @}) is refused at once. Failures
  are written to the audit log (\code{login\_failed}, with at most 64 characters of the name).
\item \textbf{Sessions.} A successful sign-in creates a 256-bit random token, and only its
  SHA-256 hash is stored. The cookie (\code{medorama\_<port>}) is \code{HttpOnly},
  \code{SameSite=Strict} and \code{Path=/}, and it is \code{Secure} when
  \code{secure\_cookies = true}. A session ends after \code{session\_hours}, or when it has not
  been used for \code{session\_idle\_minutes}. The page's automatic checks for new data do not
  count as use, so an unattended page signs out.
\item \textbf{Privacy Mode.} The session keeps its Privacy Mode choice (in its row in the
  database): a new sign-in starts in the user's default mode.
\item \textbf{Signing out.} \ui{Sign out} ends the session, and the page loads again, empty, for
  whoever signs in next. \ui{Sign out everywhere} ends all of
  the user's sessions on every device, and an administrator can do the same for a user
  (\ui{Admin} $\rightarrow$ \ui{Users} $\rightarrow$ \ui{Sign out}). A new password or a
  disabled account also ends the user's sessions.
\item \textbf{Changing a password.} Users change their own password from the user menu.
  Administrators reset other users' passwords in \ui{Admin} $\rightarrow$ \ui{Users}, or with
  \cmd{medorama user set NAME -{}-password}.
\end{itemize}

\section{Sign-in through the institution's proxy}
\label{sec:sign-in-proxy}

With \code{sign\_in = "proxy"}, an \emph{authenticating reverse proxy} stands in front of
Medorama. The proxy signs the user in with the institution's method, then passes each request to
Medorama with a header that names the user. Proxies that do this include oauth2-proxy, Apache
httpd with mod\_auth\_openidc, mod\_shib (Shibboleth) or mod\_auth\_gssapi (Kerberos), and nginx
with \code{auth\_request}. \Cref{tab:sign-in-proxy} shows how Medorama then answers
(\file{signin.rs}, \code{identify}).

\begin{xltabular}{\linewidth}{@{}P{5.4cm}L@{}}
\caption{How Medorama answers in proxy mode.}
\label{tab:sign-in-proxy}\\
\toprule
\tabhead Situation & \tabhead Result \\
\midrule
\endfirsthead
\toprule
\tabhead Situation & \tabhead Result \\
\midrule
\endhead
\bottomrule
\endfoot
The request comes from an address in \code{proxy\_addresses}, and the header names an enabled
  account & The user is signed in. Names are compared without regard to case:
  \code{JDoe@Hospital.org} matches the account \code{jdoe@hospital.org}. \\
The header names no account, a disabled account, or two accounts that differ only in case &
  HTTP 403 with the reason: ``your institution's sign-in names you \emph{name}, and Medorama has
  no account of that name: ask an administrator''. The page shows this explanation. \\
The header is missing or empty & Not signed in (HTTP 401). The page says that the institution's
  sign-in has not signed the user in. \\
The request comes from any other address & The header is ignored, and the request is not signed
  in (HTTP 401). An IPv4 address and IPv6's form of it (\code{::ffff:10.0.0.5}) are the same
  address. \\
The header is given more than once & Not signed in (HTTP 401): a copy sent by the client was
  not removed. \\
\code{POST api/v1/login} (the password form) & Refused (HTTP 400), because the institution's
  sign-in names the user. \\
A password change (\code{POST api/v1/me/password}) & Refused (HTTP 400). The page hides
  \ui{Change password} and \ui{Sign out everywhere}. \\
\ui{Sign out} & Medorama ends the visit and forgets its Privacy Mode choice, and the page goes
  to \code{proxy\_logout\_url}, the proxy's own sign-out. \\
Privacy Mode & The choice made with the switch lasts the visit (at most \code{session\_hours}),
  kept in the memory of the Medorama server that was asked; a new visit starts in the user's
  default mode. With several servers behind the proxy, another server answers in the user's
  default mode. \\
\end{xltabular}

\paragraph{The audit log.} Medorama writes one \code{login} entry per visit, with the detail
\code{proxy}, rather than one per request. A visit starts when the user has not been seen for
\code{session\_idle\_minutes} (at least one minute). An unknown name produces one
\code{login\_failed} entry per visit, with the name and the detail \code{proxy: no such account}.

\paragraph{Sessions.} In proxy mode, Medorama keeps no session of its own: the proxy's session is
the user's session, and its lifetime and idle timeout are set at the proxy or at the identity
provider. To stop a user's access at once, disable the account in Medorama, which takes effect
at the user's next request, or end the user's session at the identity provider.

\section{Security requirements for proxy sign-in}
\label{sec:sign-in-requirements}

In proxy mode, Medorama trusts the proxy to say who the user is. That trust is safe only when
the following conditions hold.
\begin{enumerate}
\item \textbf{Medorama is reachable only through the proxy.} Keep \code{host = "127.0.0.1"} and
  run the proxy on the same host. If the proxy runs on another host, bind Medorama to an
  internal address and allow connections to its port only from the proxy, with a host firewall
  or a network security group. Anyone who can connect to Medorama's port from an address in
  \code{proxy\_addresses} can name any user.
\item \textbf{Medorama's host has no untrusted accounts.} With the default
  \code{proxy\_addresses}, any program on the host can connect to 127.0.0.1 and send the header.
  Run Medorama on a host whose local accounts are trusted, such as a dedicated server or virtual
  machine, as the deployment checklist recommends (\cref{sec:security-checklist}).
\item \textbf{The proxy removes any incoming copy of the header.} A user must not be able to
  send \code{X-Remote-User} and have the proxy pass it on. Most proxies replace the header they
  set, and the examples in \cref{sec:sign-in-examples} also remove it explicitly. Check this
  after setup (step 6 of \cref{sec:sign-in-setup}).
\item \textbf{TLS terminates at the proxy.} The proxy serves HTTPS. Set
  \code{secure\_cookies = true} in \file{medorama.toml}.
\item \textbf{The address list is right.} \code{proxy\_addresses} names the address that the
  proxy's connections come from, as Medorama sees it. When Medorama listens on an IPv6 socket, an
  IPv4 proxy may appear as an IPv4-mapped address (\code{::ffff:10.0.0.5}); Medorama compares it
  with the IPv4 address listed (\code{10.0.0.5}).
\end{enumerate}

Cross-site request forgery remains blocked in proxy mode: every request that changes something
must carry the header \code{X-Medorama: 1}, which a page on another site cannot add
(\cref{sec:security-web}).

\section{Set up proxy sign-in}
\label{sec:sign-in-setup}

\begin{enumerate}
\item \textbf{Choose the account names.} Decide which value the proxy will send. It is usually
  the user's e-mail address, user principal name or eduPersonPrincipalName, such as
  \code{jdoe@hospital.org}. Medorama's account names may contain letters, digits and the
  characters \code{.}, \code{\_}, \code{-} and \code{@}, up to 64 characters, so a value such as
  \code{HOSPITAL\textbackslash jdoe} must be mapped by the proxy to an allowed form.
\item \textbf{Create the accounts} under those names, including at least one administrator
  (\cref{sec:sign-in-accounts}).
\item \textbf{Configure the proxy} to authenticate users, remove any incoming copy of the
  header, set the header to the user's name, and pass requests to Medorama
  (\cref{sec:sign-in-examples}).
\item \textbf{Change \file{medorama.toml}.} For a proxy on the same host, the settings are:
\begin{lstlisting}[language=TOML]
host = "127.0.0.1"
secure_cookies = true
sign_in = "proxy"
proxy_user_header = "X-Remote-User"
proxy_addresses = ["127.0.0.1", "::1"]
proxy_logout_url = "/oauth2/sign_out"
\end{lstlisting}
\item \textbf{Restart Medorama}, and confirm the method with the following command on the
  host. The answer names the method (\code{"method":"proxy"}) and the sign-out address.
\begin{lstlisting}[style=shell]
curl -s http://127.0.0.1:8090/api/v1/sign-in
\end{lstlisting}
\item \textbf{Check that the header cannot be forged.} From a computer other than the server,
  send the header that Medorama reads (here \code{X-Remote-User}) through the proxy without
  signing in. The proxy must redirect to its sign-in or refuse the request; it must not return a
  user:
\begin{lstlisting}[style=shell]
curl -s -i -H 'X-Remote-User: admin' https://medorama.hospital.org/api/v1/me
\end{lstlisting}
  Then try Medorama's port directly from that computer (\code{http://<server>:8090/}): the
  connection must fail.
\item \textbf{Sign in through the browser} as a user, and check \ui{Admin} $\rightarrow$
  \ui{Audit} for the \code{login} entry.
\end{enumerate}

\section{Example proxy configurations}
\label{sec:sign-in-examples}

The configurations below are short examples to adapt, not complete deployments. They show the
parts that matter to Medorama: authenticating the user, setting the header, and passing requests
to Medorama on 127.0.0.1:8090. Names in angle brackets, host names and file paths are
placeholders. Register Medorama with the identity provider as the proxy's documentation
describes, keep secrets out of configuration files where the proxy allows it, and test each
configuration with step 6 of \cref{sec:sign-in-setup}. The repository's \file{docs/sign-in.md}
gives the same examples with a few more variants.

\subsection{oauth2-proxy with an OpenID Connect provider}
\label{sec:sign-in-oauth2-proxy}

oauth2-proxy signs users in with an OpenID Connect provider and passes requests on to Medorama.
Here nginx terminates TLS in front of it. The issuer URL depends on the provider
(\cref{tab:sign-in-issuers}).

\begin{xltabular}{\linewidth}{@{}P{4.6cm}L@{}}
\caption{The issuer URL of common OpenID Connect providers.}
\label{tab:sign-in-issuers}\\
\toprule
\tabhead Provider & \tabhead \code{oidc\_issuer\_url} \\
\midrule
\endfirsthead
\toprule
\tabhead Provider & \tabhead \code{oidc\_issuer\_url} \\
\midrule
\endhead
\bottomrule
\endfoot
Microsoft Entra ID (Azure AD) & \code{https://login.microsoftonline.com/<tenant-id>/v2.0} \\
Okta & \code{https://<your-okta-domain>/oauth2/default} \\
Google & \code{https://accounts.google.com} \\
\end{xltabular}

The oauth2-proxy configuration file:
\begin{lstlisting}[language=TOML]
# oauth2-proxy.cfg (an example to adapt)
http_address  = "127.0.0.1:4180"
reverse_proxy = true                       # behind nginx, which terminates TLS
upstreams     = ["http://127.0.0.1:8090/"] # Medorama

provider        = "oidc"
oidc_issuer_url = "https://login.microsoftonline.com/<tenant-id>/v2.0"
client_id       = "<application (client) ID>"
redirect_url    = "https://medorama.hospital.org/oauth2/callback"
scope           = "openid email profile"
email_domains   = ["hospital.org"]
# client_secret and cookie_secret: in the environment (OAUTH2_PROXY_CLIENT_SECRET,
# OAUTH2_PROXY_COOKIE_SECRET), rather than in this file

cookie_secure        = true
pass_user_headers    = true                # send X-Forwarded-User and X-Forwarded-Email
prefer_email_to_user = true                # X-Forwarded-User holds the e-mail address
\end{lstlisting}

The nginx server in front of it removes any copy of the user headers that a client sends:
\begin{lstlisting}[style=text]
server {
    listen 443 ssl;
    server_name medorama.hospital.org;
    ssl_certificate     /etc/ssl/certs/medorama.pem;
    ssl_certificate_key /etc/ssl/private/medorama.key;
    client_max_body_size 1g;               # Medorama accepts imports up to 1 GiB

    location / {
        proxy_pass http://127.0.0.1:4180;  # oauth2-proxy
        proxy_set_header Host $host;
        proxy_set_header X-Forwarded-Proto $scheme;
        proxy_set_header X-Forwarded-For $remote_addr;
        proxy_set_header X-Forwarded-User "";   # never a client's own copy
        proxy_set_header X-Forwarded-Email "";
    }
}
\end{lstlisting}

The matching settings in \file{medorama.toml} are:
\begin{lstlisting}[language=TOML]
sign_in = "proxy"
proxy_user_header = "X-Forwarded-User"
proxy_addresses = ["127.0.0.1"]
proxy_logout_url = "/oauth2/sign_out"
secure_cookies = true
\end{lstlisting}

The address \code{/oauth2/sign\_out} ends the oauth2-proxy session. To end the session at the
identity provider too, and to choose another claim than \code{email} for the user's name, see
oauth2-proxy's documentation.

\paragraph{nginx with auth\_request.} nginx can also ask oauth2-proxy whether each request is
signed in (\code{auth\_request}) and pass requests to Medorama itself. oauth2-proxy then runs with
\code{upstreams = ["static://202"]} and \code{set\_xauthrequest = true}, so that it answers
nginx's checks with the user's address in \code{X-Auth-Request-Email}. In nginx, the location
that passes requests to Medorama copies that address into \code{X-Remote-User}, which replaces
any copy the client sent (the \code{/oauth2/} locations are as in oauth2-proxy's documentation):
\begin{lstlisting}[style=text]
location / {
    auth_request /oauth2/auth;
    error_page 401 =403 /oauth2/sign_in;
    auth_request_set $user $upstream_http_x_auth_request_email;
    proxy_set_header X-Remote-User $user;   # replaces any copy the client sent
    proxy_pass http://127.0.0.1:8090;       # Medorama
}
\end{lstlisting}
Keep the default header in \file{medorama.toml}, with \code{proxy\_addresses = ["127.0.0.1"]}
and \code{proxy\_logout\_url = "/oauth2/sign\_out"}.

\subsection{Apache httpd with mod\_auth\_openidc}
\label{sec:sign-in-openidc}

mod\_auth\_openidc adds OpenID Connect sign-in to Apache httpd. It sets the user's name
(\code{REMOTE\_USER}) from the claim named by \code{OIDCRemoteUserClaim}, and mod\_headers passes
it to Medorama after removing any copy the client sent:
\begin{lstlisting}[style=text]
<VirtualHost *:443>
    ServerName medorama.hospital.org
    SSLEngine on
    SSLCertificateFile    /etc/pki/tls/certs/medorama.pem
    SSLCertificateKeyFile /etc/pki/tls/private/medorama.key

    OIDCProviderMetadataURL https://<provider>/.well-known/openid-configuration
    OIDCClientID            <application (client) ID>
    OIDCClientSecret        <client secret>
    OIDCRedirectURI         https://medorama.hospital.org/redirect_uri
    OIDCCryptoPassphrase    <a long random string>
    OIDCScope               "openid email profile"
    OIDCRemoteUserClaim     email

    <Location />
        AuthType openid-connect
        Require valid-user
    </Location>

    RequestHeader unset X-Remote-User early
    RequestHeader set X-Remote-User expr=%{REMOTE_USER}
    ProxyPass        / http://127.0.0.1:8090/
    ProxyPassReverse / http://127.0.0.1:8090/
</VirtualHost>
\end{lstlisting}
The \code{expr=} form of \code{RequestHeader} needs Apache httpd 2.4.10 or later. If the
provider sends no \code{email} claim, name another one, such as \code{preferred\_username}. For
sign-out, mod\_auth\_openidc accepts a \code{logout} parameter on its redirect URI, for example
\code{proxy\_logout\_url = "/redirect\_uri?logout=https\%3A\%2F\%2Fmedorama.hospital.org\%2F"}.

\subsection{Shibboleth (SAML), Kerberos, and VPNs}
\label{sec:sign-in-other-proxies}

\paragraph{Shibboleth.} The Shibboleth Service Provider (mod\_shib) signs users in through a SAML
identity provider, such as an institution's or a research federation's. The attribute that
becomes \code{REMOTE\_USER} is chosen in \file{shibboleth2.xml} (for example \code{eppn}). The
virtual host is the same as for mod\_auth\_openidc, with \code{AuthType shibboleth},
\code{ShibRequestSetting requireSession 1} and \code{Require shib-session} in its
\code{<Location />}, and \code{ProxyPass /Shibboleth.sso !} before the other \code{ProxyPass}
line, so that the Service Provider's own handlers stay in Apache httpd. Keep mod\_shib's default
of passing attributes as environment variables rather than headers, and set
\code{proxy\_logout\_url = "/Shibboleth.sso/Logout"}.

\paragraph{Kerberos.} On a network where users' computers hold Kerberos tickets (for example an
Active Directory domain), mod\_auth\_gssapi signs them in without a password prompt:
\code{AuthType GSSAPI}, \code{GssapiCredStore keytab:/etc/httpd/medorama.keytab},
\code{GssapiLocalName On} (so that \code{jdoe@HOSPITAL.ORG} becomes \code{jdoe}) and
\code{Require valid-user}, with the same \code{RequestHeader} and \code{ProxyPass} lines in a TLS
virtual host. Without \code{GssapiLocalName}, the name is the full principal, which is also a
valid Medorama account name. Kerberos has no sign-out, so leave \code{proxy\_logout\_url} empty.

\paragraph{A VPN or an access gateway.} A VPN restricts who can reach Medorama, but most VPNs do
not tell a web application who the user is. Behind a VPN, either keep
\code{sign\_in = "password"}, so that users sign in to Medorama inside the VPN, or add one of the
proxies above. An access gateway or identity-aware proxy that adds a header naming the user works
like the proxies above, provided that the header holds exactly the account name. Medorama
compares the whole value, so a value with a prefix (such as a provider's name and a colon) must
be rewritten by a proxy step first. List the gateway's addresses in \code{proxy\_addresses}.

\section{Accounts, roles and the PHI permission}
\label{sec:sign-in-accounts}

Accounts live in Medorama with either method. In proxy mode, create each account under the exact
name the proxy sends; case does not matter.
\begin{itemize}
\item \textbf{The first administrator} is created by \cmd{medorama init}. In proxy mode, give it
  the name the proxy will send, for example
  \cmd{medorama init -{}-results DIR -{}-admin jdoe@hospital.org}.
\item \textbf{More accounts} are created in \ui{Admin} $\rightarrow$ \ui{Users}, or on the server
  with \cmd{medorama user add}. In proxy mode an account needs no password. The following
  command adds an analyst with the PHI permission who signs in through the institution's
  sign-in:
\begin{lstlisting}[style=shell]
medorama user add jdoe@hospital.org --role analyst --phi --display-name "Jane Doe" --no-password
\end{lstlisting}
  With \opt{-{}-no-password}, the account receives a random password that nobody knows, so it can
  never sign in with a password, even if the catalog is later switched back to password
  sign-in (an administrator then sets a password with \cmd{medorama user set NAME
  -{}-password}). With proxy sign-in, \ui{Admin} $\rightarrow$ \ui{Users} hides the password
  field of a new account and the \ui{Reset password} button, and \code{POST users} may leave the
  password out. Without \opt{-{}-no-password}, the command asks for a password of at least 10
  characters.
\item \textbf{Changes} are made in \ui{Admin} $\rightarrow$ \ui{Users}, or with
  \cmd{medorama user set NAME} (\opt{-{}-role}, \opt{-{}-phi true|false},
  \opt{-{}-disabled true|false}, \opt{-{}-display-name}). \cmd{medorama user list} lists the
  accounts, and \cmd{medorama user remove NAME} deletes one. Each change is written to the audit
  log.
\end{itemize}

The roles and the PHI permission are described in \cref{chap:users-privacy}. Medorama does not
read groups or roles from the identity provider: an administrator assigns them in Medorama.

\section{Adding another sign-in method in code}
\label{sec:sign-in-code}

Most institutions need no code: a proxy in front of Medorama connects it to any sign-in that the
proxy supports. If a method must be built into Medorama instead (for example, verifying a signed
token that a gateway adds to each request), the change is confined to the sign-in code.
\begin{enumerate}
\item \textbf{Add a variant to \code{signin::Method}} in \file{medorama/src/signin.rs}, holding
  the method's settings, and give it a name in \code{Method::name}.
\item \textbf{Add its arm to \code{signin::identify}.} From the request's headers, its peer
  address and the database, return \code{Who::User(user, choice)} for a signed-in user (with
  the session's Privacy Mode choice, \code{None} if none was made: kept by
  \code{signin::set\_privacy}),
  \code{Who::Nobody} when the request is not signed in, or \code{Who::NoAccount(name)} when the
  method names someone without an account. Look the account up with
  \code{users::get\_by\_login}.
\item \textbf{If users sign in through Medorama's page} (a form, or a redirect that returns to
  Medorama), handle it in \code{signin::login}, and in \code{signin::logout} if signing out needs
  more than forgetting the Privacy Mode choice. If the method needs its own sign-in screen, the
  page's \code{showLogin} in \file{medorama/web/js/main.js} changes too.
\item \textbf{Read the method's settings} in \file{config.rs}: add its value to
  \code{sign\_in\_method} and its keys to \code{Config::load}, so that a mistake stops Medorama at
  start.
\item \textbf{Test it.} Add a test to \file{medorama/tests/signin.rs} that signs in, is refused
  from the wrong place, and checks the audit log, and extend the unit test
  \code{sign\_in\_settings} of \file{config.rs} with the new settings and their errors.
\end{enumerate}
Nothing else changes: every handler receives the signed-in user as \code{server::Auth(User)},
and roles, the PHI permission, Privacy Mode and the privacy gate apply as before.

\section{The viewers' own sign-in}
\label{sec:sign-in-viewers}

The viewers can also run alone, without Medorama, on a researcher's workstation or a server
(\cref{chap:viewers-alone}).
\begin{itemize}
\item \textbf{VitalStream} (\cmd{vitalstream-py}, \cmd{vitalstream-rs}, \cmd{vitalstream-cpp})
  listens on 127.0.0.1 by default (ports 8050, 8060 and 8070). \cmd{vitalstream-rs
  -{}-set-password}, or the same option of the other two viewers, sets an optional password that
  the page then asks for. It is stored as a salted scrypt hash, with a random signing key, in
  \file{\textasciitilde/.config/vitalstream/auth.json}, readable only by its owner, and serves
  all three viewers; a sign-in lasts 30 days, and signing out ends it on the server too. After
  wrong passwords from one address, the next one is checked only after 1, 3, 7, 15 \dots{} up to
  30\,s. Set a password on any computer that other people use.
\item \textbf{A build without password support.} \cmd{viewers/vitalstream/install.sh
  -{}-no-password} builds a viewer with no password code at all (for the Rust viewer
  \cmd{cargo build -{}-no-default-features -{}-features wfdb}, for the C++ viewer the CMake option
  \code{-DVS\_PASSWORD=OFF}). Such a build refuses to run where a password is required, as under
  Medorama.
\item \textbf{A root folder fixed when the viewer is built.} The viewer never opens anything
  outside its root: \cmd{install.sh -{}-root DIR} fixes it (by default \file{\textasciitilde}, the
  home folder of whoever runs the viewer). The \code{data\_roots} setting can narrow the root but
  never widen it, and symbolic links are resolved before the check.
\item \textbf{The ECG and PACS viewers} are web pages. Opened from disk, they read only the files
  the user chooses in the browser and need no sign-in. Served with \cmd{muse-py serve} or
  \cmd{dicom-py serve}, they list and open files under a root (\opt{-{}-root}, inside the root
  fixed at installation with \code{VIEWER\_ROOT}), listen on 127.0.0.1 (ports 8080 and 8081), and
  accept the same optional password (\opt{-{}-set-password}; a package installed with
  \code{VIEWER\_PASSWORD=0} has no password support). For \cmd{muse-py serve} and
  \cmd{dicom-py serve}, a sign-in lasts 30 days, until the server stops, or until a browser
  signs out, which ends every sign-in; after a wrong password the next try, from any browser,
  waits 1\,s, then 2, 4 \dots{} up to 60\,s.
\item \textbf{Host names.} Every viewer answers only requests that name it by a loopback name,
  an IP address or its own host, which defeats DNS rebinding; VitalStream's
  \code{allowed\_hosts} adds the names it is reached by (\cref{sec:viewer-api-checks}).
\end{itemize}
To reach a viewer on a server from another computer, use an SSH tunnel, for example
\cmd{ssh -N -L 8060:127.0.0.1:8060 user@server}.

\textbf{Inside Medorama}, the viewers sit behind Medorama's sign-in. The ECG and PACS pages are
served by Medorama and read a patient's files through its API, like every other request. The
VitalStream viewer runs as a separate program on a private port of 127.0.0.1, with a password
file (\file{<home>/viewer/auth.json}, mode 0600) whose signing key only Medorama knows.
Medorama's gateway signs every request it passes on, users never see the viewer's sign-in, and
the viewer's port is of no use to other accounts on the host (\cref{sec:architecture-gateway}).
